% !TeX root = ../notes.tex
% Proposed prose version aligned to 08-autonomise-plan.tex.

\chapter{From automated synthesis to self-driving chemistry}

\begin{notebox}
  \keypoint{What do you need to learn from this chapter?}
  Automation, programmability, and autonomy describe different relationships between chemical procedures, laboratory hardware, measurements, and decisions.
  A self-driving laboratory closes a tested loop in which reliable measurements influence subsequent experimental actions.
\end{notebox}

\section{Automation, programmability, and autonomy}

Automation executes a predefined sequence of operations.
It can improve reproducibility, throughput, safety, and access to conditions that are difficult to maintain manually.
The scientific decisions may nevertheless have been made completely before the run begins.

Programmability represents laboratory procedures in a digital form that can be stored, edited, versioned, and executed on compatible hardware.
A programmable system separates the logical procedure from a particular sequence of manual actions.
It does not become autonomous merely because the procedure is written as code.

Autonomy requires a measurement to influence what the system does next.
The decision can use a fitted model, an acquisition function, a heuristic rule, or a threshold.
The essential feature is feedback from observed outcome to subsequent action.

\begin{conceptbox}
  \keypoint{Automation executes.}
  \keypoint{Programmability represents and transfers.}
  \keypoint{Autonomy measures and decides.}
\end{conceptbox}

\section{Automated peptide synthesis}

Merrifield and Stewart's automated peptide synthesiser mechanised the repeated cycle required for solid-phase peptide synthesis.\mcite{Merrifield1965}
Deprotection, washing, coupling, and washing can be performed in a predefined order for each residue in the sequence.
The achievement is important because reliable repetition converts a demanding manual procedure into a scalable synthetic operation.

The synthesiser does not decide which amino acid should be coupled after measuring the growing peptide.
The complete sequence and operating logic are supplied before synthesis.
It is therefore a clear historical example of automation without autonomy.

\section{Chemputer and programmable synthesis}

The Chemputer developed by Steiner and co-workers represents organic synthesis through programmable unit operations.\mcite{Steiner2019}
A written procedure is translated into chemical code, combined with a graph describing the physical platform, and compiled into commands for pumps, valves, heaters, separators, filters, and evaporators.
The abstraction allows reaction, work-up, isolation, and purification to be represented in one executable workflow.

The reported syntheses of diphenhydramine, rufinamide, and sildenafil demonstrate that multistep procedures can be transferred into machine execution.
The platform's central advance is programmability.
The target and sequence are specified by the chemist rather than selected adaptively from analytical feedback.

\section{From automated experimentation to a closed loop}

A closed experimental loop requires more than a robot and an optimisation algorithm.
The platform must prepare and execute experiments, acquire analytical data, convert instrument output into a quality-controlled response, update its state of knowledge, and select a subsequent action.
Failures at any interface can break the scientific meaning of the loop even when every device continues to operate.

Chapter~6 supplied strategies for choosing experiments from objectives, predictions, and uncertainty.
Chapter~7 supplied execution, analysis, provenance, and failure states.
Autonomy connects these capabilities so that a validated measurement causes a new experimental decision.

\section{Cooper: a mobile robotic chemist}

Burger and co-workers used a mobile robot to operate instruments distributed through a conventional laboratory.\mcite{Burger2020}
The robot prepared photocatalyst formulations, transported samples, initiated experiments, and returned measurements to a batched Bayesian search.
The campaign performed 688 experiments over eight days in a ten-variable space and identified mixtures substantially more active than the initial formulations.

Mobility solves an orchestration problem because one robotic agent can connect instruments that were not designed as one enclosed platform.
Bayesian selection supplies the autonomous element because completed measurements influence later batches.
The case therefore combines physical mobility, automated execution, and adaptive decision-making.

\section{RoboChem}

RoboChem integrates flow photochemistry, phase tracking, inline NMR, and Bayesian optimisation in a closed loop.\mcite{Slattery2024}
The platform can optimise continuous and categorical variables against objectives such as yield and productivity.
It also links optimisation with process intensification and scale-up.

The distinction between these tasks should remain explicit.
Optimisation identifies favourable conditions within the defined search space.
Intensification changes how efficiently material is produced, while scale-up tests whether the selected conditions transfer to a larger production objective.
The platform is autonomous because measured reaction outcomes update the decision about the next conditions.

\section{Autonomous exploratory synthesis}

Dai and co-workers connected an automated synthesiser, UPLC-MS, and benchtop NMR using mobile robots and an application-agnostic heuristic planner.\mcite{Dai2024}
A reaction advances only when orthogonal analytical measurements pass predefined, application-specific thresholds.
The system can then request replication, scale-up, diversification, or a functional host--guest assay.

These workflows are autonomous because measurements determine subsequent actions.
They are not closed-loop optimisation processes because they do not repeatedly update a scalar response model to improve an objective.
The distinction is productive rather than limiting.
Exploratory chemistry can require branching decisions among identity, reproducibility, scale-up, and function rather than movement towards one numerical optimum.

The study also exposes the limits of automated context.
An unexpected cyclisation was not recognised by the initial UPLC-MS rule and required later human interpretation.
Orthogonal measurements reduce ambiguity, but an autonomous planner can reason only with the signals, thresholds, and chemical representations available to it.

\begin{warningbox}
  \keypoint{Autonomous does not necessarily mean optimising.}
  A system can make consequential experimental decisions through rules and analytical thresholds without fitting and updating an optimisation model.
\end{warningbox}

\section{Robustness, error handling, and recovery}

An autonomous platform must distinguish a chemical negative result from a hardware or analytical failure.
A blocked line, missed injection, exhausted reagent, detector drift, failed integration, and genuinely inactive reaction require different responses.
Quality-control rules should determine whether a result is accepted, repeated, diagnosed, or escalated.

Safe recovery includes instrument watchdogs, calibration checks, inventory tracking, bounded retries, and a defined shutdown state.
Every rejected or repeated experiment should retain a traceable reason.
Without this record, the decision history cannot be audited and a later model may treat fault-generated values as chemistry.

\section{Human-in-the-loop and responsible autonomy}

Humans define the chemical problem, objective, constraints, permitted actions, and evidence required for success.
The autonomous system chooses only within that delegated decision space.
Chemical judgement remains necessary when objectives conflict, measurements are ambiguous, or the system encounters chemistry outside its representation.

Responsible deployment also requires considering resource use and safety.
An objective that rewards yield alone can direct a platform towards expensive, hazardous, or wasteful conditions unless those consequences are represented as objectives or constraints.
Explainability should be sufficient to reconstruct why an experiment was selected and which evidence caused the next action.

\section{Closing the course}

The course began by translating molecules into representations that computers can compare.
It then used descriptors to map chemical space, models to connect representation with response, selection strategies to choose experiments, and automated platforms to execute and measure them.
Autonomy closes the loop:
\begin{equation}
  \text{select}\longrightarrow\text{execute}\longrightarrow\text{measure}\longrightarrow\text{update}\longrightarrow\text{select}.
\end{equation}

Every stage inherits limitations from the earlier stages.
A poor representation hides relevant chemistry, a biased dataset distorts the model, an inappropriate objective rewards the wrong outcome, and an unreliable measurement corrupts the feedback signal.
A self-driving laboratory is therefore not one instrument or one algorithm.
It is a validated relationship between a chemical objective, executable operations, trustworthy measurements, explicit decision logic, and safe recovery.

\section{Summary}

Automation, programmability, and autonomy should be distinguished by what is executed, represented, measured, and decided.
Automated peptide synthesis demonstrates predefined repetitive execution.
The Chemputer demonstrates transferable digital procedures.
The mobile robotic chemist and RoboChem close optimisation loops through measured feedback.
The 2024 exploratory-synthesis platform demonstrates autonomous branching without claiming closed-loop optimisation.
Human judgement remains responsible for objectives, constraints, interpretation, and intervention.
